\paragraph{Analytic radii of the imported leading profiles.}
The leading fields of the hierarchy are the source profiles from
Appendix~B of \cite{openai2026}, with $Y=\Lambda X$.
Proposition~B.2 constructs these fields on $0\le Y\le4.1$
and a common neighborhood of $[-1,1]$ in $\eta$, once
admissible $\Lambda,C$ and axis data have been chosen. Its
proof selects a radius $S$ with $4.1<S<20$ and a smaller
positive complex $\eta$-neighborhood. This proves qualitative
regularity for fixed admissible parameters; it does not supply
explicit numerical values of $\Lambda,C,S,\delta$, nor a
certified complex supremum for one frozen realization.

\begin{lemma}[Source-radius transfer into the doubled norm]
\label{lem:source-radius-transfer}
Suppose $F(Y,\eta)$ is holomorphic on $|Y|<S$ and in
the $\delta$-neighborhood of $[-1,1]$, with $|F|\le M$.
Then $F_0(X,\eta)=F(\Lambda X,\eta)$ obeys
$\|F_0\|_{\mathcal B}\le M$ provided
$R\le S/(8\Lambda)$ and $\rho\le\delta/8$.
For any fixed integers $p,q\ge0$,
\begin{equation}\label{eq:source-Cauchy-bound}
\|D_X^p\partial_\eta^qF_0\|_{\mathcal B}
\le M\delta^{-q}\sup_{a,b\ge0}
\frac{a^p(a+1)^2(b+1)^2}{8^{a+b}}
\frac{(b+q)!}{b!}<\infty.
\end{equation}
\end{lemma}
\begin{proof}
Cauchy's inequalities give
$|\partial_\eta^{b+q}F_a|\le MS^{-a}(b+q)!\delta^{-b-q}$;
substitution of $Y=\Lambda X$ supplies $\Lambda^a$.
Divide by (37) and use $\binom{2a+b}{b}^{-1}\le1$.
The radial and parameter radius ratios are at most $1/8$.
For $p=q=0$, $(a+1)^2/8^a\le1$ and
$(b+1)^2/8^b\le1$. General fixed derivatives introduce
only polynomial factors absorbed by the geometric terms.
\end{proof}
The source's $U_0,\varphi_0$ and their finite derivatives
therefore have finite multiplier budgets after shrinking the
working radii. The incompressibility formula for $V_0/X$
uses only radial averages, rational factors with $L\ne0$,
and finite parameter derivatives. These operations preserve
holomorphy on a smaller common domain. This is an
\emph{existential} certification, not a numerical enclosure
of any fixed source profile.

\paragraph{The complete current-order swirl-pressure block.}
The two ordered terms $(0,n),(n,0)$ of (27) contribute
\begin{equation}\label{eq:current-pressure-identity}
\Pi_n^{\mathrm{cur}}(F)
=\frac2{C^2}\int_0^X\varphi_0(x,\eta)F(x,\eta)\,dx,
\qquad F=\varphi_n\in I_n.
\end{equation}
They belong in $K_n^{U\varphi}$, not in lower-order forcing.
\begin{lemma}[Quadratic radial gain in current pressure]
\label{lem:current-pressure-quadratic}
For $F\in I_n(\mathcal B)$,
\begin{align}
\left\|\int_0^X F\,dx\right\|_{\mathcal P}
&\le\frac{8R\rho}{(n+1)^2}\|F\|_{\mathcal B},\label{eq:axis-primitive-n2}\\
\|\Pi_n^{\mathrm{cur}}(F)\|_{\mathcal P}
&\le\frac{1024R\rho\|\varphi_0\|_{\mathcal B}}
{|C|^2(n+1)^2}\|F\|_{\mathcal B},\label{eq:current-pressure-n2}\\
\|K_n^{U\varphi}F\|_{\mathcal B}
&\le\frac{1024R\rho C_E\|\varphi_0\|_{\mathcal B}}
{|C|^2(n+1)^2}\|F\|_{\mathcal B},\label{eq:current-block-n2}
\end{align}
where $C_E$ is the pressure-to-velocity constant in
\eqref{eq:P-to-U}.
\end{lemma}
\begin{proof}
Factorial cancellation gives the exact weight ratio
\[
\frac{w_{a,b}}{(a+1)w_{a+1,b+1}}
=R\rho\frac{(a+2)^2(b+2)^2(2a+1)(2a+2)}
{(a+1)^3(b+1)^2(2a+b+1)(2a+b+2)(2a+b+3)}.
\]
Its maximum at fixed $a$ occurs for $b=0$, giving
$4R\rho(a+2)^2/[(a+1)^3(2a+3)]$.
The identity
$2(a+1)(2a+3)-(a+2)^2=3a^2+6a+2\ge0$
bounds this by $8R\rho/(a+1)^2$.
Since $a\ge n$ this proves the first inequality.
The Banach-algebra constant $64$, and the factor two
from the two ordered source pairs, prove the second.
Theorem~\ref{thm:pressure-to-velocity} proves the third.
\end{proof}

\paragraph{Conditional simultaneous pressure convergence.}
The preceding estimate replaces the earlier current-pressure
bound without double-counting $\varphi_0\varphi_n$.
If the three other current-block estimates
\eqref{eq:swirlbound}, \eqref{eq:phiuB} and \eqref{eq:uuB}
hold with a finite combined constant $C_{\rm other}$,
then
\begin{equation}\label{eq:current-K-updated}
\|K_n\|\le\frac{C_{\rm other}}{n+1}
+\frac{1024R\rho C_E\|\varphi_0\|}
{|C|^2(n+1)^2}.
\end{equation}
The Neumann-series inverse has norm at most two for
$n$ larger than twice the sum of the numerator constants.

Write $B_n=\|\varphi_n\|+\|U_n\|$, and suppose
$B_n\le MT^n/(n+1)^2$ for $n\ge1$ with $M,T\ge1$.
Then the complete shifted pressure ledger
\eqref{eq:complete-two-space-P-ledger},
after separating the current ordered pairs, implies
\begin{equation}\label{eq:conditional-pressure-geometric}
\|\Pi_n\|_{\mathcal P}
\le\frac{\Gamma T^n}{(n+1)^2},\quad n\ge3,
\end{equation}
where
\[
\Gamma=2C_{\rm sw}B_0M+4C_{\rm sw}M^2
+\frac{4C_LM}{T}+\frac{4C_DM}{T^2}
+\frac{4C_{\rm quad,P}M^2}{T},
\]
$B_0=\|\varphi_0\|+\|U_0\|$,
$C_L=C_{\rm P,time}+C_{\rm P,lap}+C_{\rm lead,P}$
and $C_D=C_{\rm P,del}$. The proof uses
$\sum_{i=1}^{n-1}[(i+1)^2(n-i+1)^2]^{-1}
\le4/(n+2)^2$ and the elementary adjacent-index
comparisons. The finite startup indices are absorbed by
enlarging $\Gamma$. This implication is not a proof that
the actual velocity sequence obeys the assumed bound.

\paragraph{Physical-coordinate reconstruction, conditional on convergence.}
With $\tau=q(1-\eta^2)$ and $z=q^D\eta$,
\begin{equation}\label{eq:coordinate-inverse-Jacobian}
\left.\frac\partial{\partial q}
[q-z^2q^{1-2D}]\right|_{z\ {\rm fixed}}
=1-2h\eta^2=L(\eta)>0.
\end{equation}
Thus the physical chain-rule coefficients are smooth on
compact subcores with $q>0$. If the velocity and pressure
weighted geometric estimates above hold on a common domain,
the radial valuation and coefficient weights give normal
convergence after each fixed $X,\eta$ derivative when
$T|q^{2h}|<2$ and $|X|<R/8$. Physical differentiation
then passes termwise on compact subcores. Exact recovery
of the unforced equations additionally requires the actual
coefficient hierarchy, including its startup normalization.

\paragraph{The finite-startup radius-selection gap.}
A bound $\|K_n\|\le C_K(\rho)/(n+1)$ may require a cutoff
$N_*(\rho)$. Merely knowing each finite-order coefficient
has some positive analytic radius $\rho_n$ does not
guarantee a common $\rho$ obeying
\begin{equation}\label{eq:radius-startup-fixed-point}
\rho\le\rho_{\rm lead},\qquad
\rho\le\min_{1\le n<N_*(\rho)}\rho_n.
\end{equation}
For example, the abstract choices $\rho_n=e^{-n}$ and
$N_*(\rho)=\lceil1/\rho\rceil+1$ provide positive radii
at every finite order but no $0<\rho\le1$ satisfying
the second inequality: its right-hand side is at most
$e^{-1/\rho}<\rho$. This does not refute the source
Navier--Stokes hierarchy. It shows that a common-radius
induction or noncircular quantitative initialization is
still required before the closed recurrence and
Theorem~1.1 can be claimed.
